\documentclass[10pt]{article}
\usepackage[utf8]{inputenc}
\usepackage[spanish, es-tabla]{babel}
\usepackage{amsmath}
\usepackage{amssymb}
\usepackage{siunitx}
\usepackage{csquotes}
\usepackage{tabularx}
\usepackage{verbatim}
\usepackage{dirtree}
\usepackage{dsfont}
\usepackage{biblatex}
\usepackage[table, dvipsnames]{xcolor}
\usepackage[most]{tcolorbox}
\usepackage{multirow}
\usepackage{multicol}
\usepackage{float}
\usepackage{caption}
\usepackage{hyperref}
\usepackage{ragged2e}
\usepackage{wrapfig}
\usepackage[labelformat=simple]{subfig}
\usepackage{pgfplots}
\pgfplotsset{compat=1.18} 
\usepackage{makecell}
\usepackage{tikz}
\usetikzlibrary{
    shapes.geometric,
    positioning,
    calc,
    automata,
    arrows
    }
\usepackage{epstopdf}
%\usepackage{subcaption}
\renewcommand{\thesubfigure}{}
%\captionsetup{subfigure}{labelformat=empty}
\usepackage{ifthen}
\usepackage{enumitem}
\usepackage{parskip}
\usepackage{geometry}
\usepackage{graphicx}
\usepackage{fancyhdr}
\geometry{a4paper, left=15mm, right=15mm, top=10mm, bottom=5mm}
\usepackage[T1]{fontenc}
\renewcommand{\thesubfigure}{\relax}  % Do nothing for the counter »subfigure«


%%%%%%%%%  Ambientes útiles para abreviar la escritura al dibujar los nodos en tikz

\newenvironment{RedC}{%
    \begin{center}%
    \begin{tikzpicture}[
        node distance = {5mm},
        main/.style = {
            draw,
            circle,
            minimum size = 4mm,
            inner sep = 0pt,
            very thick
        }
    ]%
}{%
    \end{tikzpicture}%
    \end{center}
}

\newenvironment{RedNC}{%
    \begin{tikzpicture}[
        node distance = {5.5mm},
        main/.style = {
            draw,
            circle,
            minimum size = 3.5mm,
            inner sep = 0pt,
            very thick
        }
    ]%
}{%
    \end{tikzpicture}
}

\newenvironment{RedPNC}{%
    \begin{tikzpicture}[
        node distance = {4mm},
        main/.style = {
            draw,
            circle,
            minimum size = 2.5mm,
            inner sep = 0pt,
            very thick
        }
    ]%
}{%
    \end{tikzpicture}
}

\newcommand*{\punto}{\node[main]}

\newcolumntype{C}{>{\centering\arraybackslash}X}


%%%%%%%%%%%%%%%%%%%%%%%%%%%%%%%%%%%%%%%%%%%%%%%%%%%%%%%%%%%%%%%%%%%%%%%%%%

\newcommand{\competidorA}{Nombre Competidor A}
\newcommand{\competidorB}{Nombre Competidor B}

%%%%%%%%%%%%%%%%%%%%%%%%%%%%%%%%%%%%%%%%%%%%%%%%%%%%%%%%%%%%%%%%%%%%%%%%%%



\begin{document}


\noindent
%encabezado, mitad izquierda
\begin{minipage}[t]{0.5\textwidth}
    %logo
    \begin{minipage}[t]{0.30\textwidth}
        \vspace{-7mm}
        \includegraphics[width=\textwidth]{./logo_PA.png}
    \end{minipage}%
    \hspace{2mm}
    %texto de inicio
    \begin{minipage}[t]{0.65\textwidth}
        {\Huge\textbf{Ficha Médica}}\\[6pt]
        {\Large\textbf{Lesiones y}}\\[5pt]
        {\Large\textbf{Traumatismos}}
    \end{minipage}
\end{minipage}%
\hfill
%encabezado, mitad derecha
\begin{minipage}[t]{0.45\textwidth}
    \raggedleft
    \vspace{-5mm}
    {\large\textbf{Liga Interna de}}\\[1pt]
    {\large\textbf{Espada Larga}}\\[1pt]
    {\large Edición 2026}\\[-0.5pt]
    {\large Praelium Australis}
\end{minipage}

\vspace{5mm}

\begin{center}
    \textcolor{gray}{\textit{Rellenar en caso de lesión o traumatismo sufrido por uno o más participantes del duelo}} \\
    \textcolor{gray}{\textit{Esta hoja deberá ser entregada a la persona afectada o el personal médico en caso de ser trasladada a un centro de asistencia}}
\end{center}

\vspace{5mm}

\begin{tabularx}{\textwidth}{X X}
    {\large\textbf{Identidad de la persona lesionada}} &  \\[20pt]
    \textbf{Nombre:}\hspace{1.5mm}\rule[0pt]{150mm}{1pt} & \\[10pt]
    \textbf{Rol en el Duelo} & \\[5pt]
    Competidor:\hspace{2mm}%
    \begin{RedPNC}
        \punto (p1) {} ;
    \end{RedPNC}%
    \hspace{4mm}%
    Juez/Árbitro:\hspace{2mm}%
    \begin{RedPNC}
        \punto (p1) {} ;
    \end{RedPNC}%
    \hspace{4mm}
    Espectador:\hspace{2mm}%
    \begin{RedPNC}
        \punto (p1) {} ;
    \end{RedPNC}
\end{tabularx}

\begin{tabularx}{\textwidth}{| X | X |}\hline
    \vspace{1mm}{\large\textbf{Zona afectada}} & \vspace{1mm}{\large\textbf{Descripción}}\\[10pt]
    \begin{itemize}
        \item[$\dagger$] Cabeza
    \end{itemize}
\end{tabularx}



\end{document}